\cvsection{Expérience professionnelle}
\begin{cventries}

  \cventry
    {Chercheur quantitatif}
    {Synerqo}
    {Montréal, Canada}
    {Mai 2025 - Aujourd'hui}
    {
      \begin{cvitems}
        \item{Développement d'un \textbf{CRM sur mesure et d'un portail conseiller pour un client du secteur des assurances}, intégrant traitement des commissions, signature électronique et tableaux de bord commerciaux, avec une optimisation des processus métier et \textbf{plusieurs dizaines de milliers de dollars d'économies}.}
        \item{Conception, développement et déploiement de \textbf{\href{https://ganesha.synerqo.com/fr}{Ganesha}}, l'assistant d'IA de Synerqo pour les entreprises de services professionnels, combinant recherche documentaire et IA générative (RAG), réponses reliées aux sources, génération de documents et processus adaptés à la documentation et aux méthodes de chaque client.}
        \item{Construction d'un \textbf{moteur de backtest événementiel sur options en Python et Numba}, intégrant écarts acheteur-vendeur, glissement, commissions, marge et assignation pour valider les stratégies avant mise en production, en complément de QuantConnect.}
        \item{Développement d'un \textbf{moteur de valorisation d'options exotiques et de produits structurés} avec Monte Carlo, solveurs d'EDP, sensibilités par \textbf{différentiation algorithmique adjointe (AAD)}, couverture automatisée et calibration Heston, SABR et SVI.}
        \item{Développement de \textbf{composants d'exécution en Rust} pour l'ingestion des données de marché, les cotations et le routage des ordres, avec mesure de la latence face à l'environnement de recherche Python.}
        \item{Conception, backtest et \textbf{mise en production d'une stratégie systématique sur options SPX}, issue du mémoire de maîtrise sur l'exposition gamma des teneurs de marché.}
        \item{\textbf{Technologies :} Python, C++, Rust, TypeScript, React/Next.js, FastAPI, PostgreSQL, QuantLib, QuantConnect.}
      \end{cvitems}
    }

  \cventry
    {Analyste quantitatif | Ingénierie financière et modélisation des risques}
    {Deloitte}
    {Montréal, Canada}
    {Octobre 2024 - Juin 2025}
    {
      \begin{cvitems}
        \item{Développement et validation de \textbf{modèles de risque de crédit (PD, LGD, EAD)} pour de grandes institutions financières canadiennes, dans les cadres IFRS 9 et Bâle III/IV.}
        \item{Pilotage de missions de validation indépendante avec analyse des hypothèses, de la calibration et du pouvoir discriminant, tests rétrospectifs, analyses de sensibilité et tests de résistance.}
        \item{Création de \textbf{tableaux de bord de suivi et de chaînes de validation automatisées en Python}, réduisant de \textbf{40\,\%} le temps de revue manuelle.}
        \item{Rédaction de rapports techniques et de documentation réglementaire, puis présentation des conclusions aux \textbf{dirigeants et comités de risque des clients}, dans le cadre des exigences du BSIF.}
        \item{\textbf{Technologies :} Python, R, SQL, modélisation statistique, suivi et validation de modèles.}
      \end{cvitems}
    }

  \cventry
    {Stagiaire développeur quantitatif | R\&D, dérivés sur actions}
    {Banque Nationale Marchés financiers}
    {Montréal, Canada}
    {Mai 2023 - Décembre 2023}
    {
      \begin{cvitems}
        \item{Réalisation de \textbf{1 M\$ d'économies annuelles} grâce à l'optimisation des stratégies et à la migration d'un fournisseur externe vers des systèmes de production internes.}
        \item{Construction d'une \textbf{plateforme complète de backtest en Python} couvrant signaux, construction de portefeuille, coûts de transaction, analyses de risque et attribution de performance.}
        \item{Développement et déploiement de stratégies de faible volatilité et de momentum \textbf{SmartBeta couvrant 800 M\$ d'actifs sous gestion}, au sein d'une division de produits structurés de 7 G\$, avec optimisation de portefeuille sous contraintes via cvxpy.}
        \item{Collaboration avec les \textbf{équipes de structuration et de négociation} pour intégrer les signaux aux outils de valorisation et de couverture, et rapports de performance et de risque destinés aux négociateurs et gestionnaires.}
        \item{\textbf{Technologies :} Python, cvxpy, SQL, terminal et API Bloomberg, AWS (EC2, S3), tests statistiques.}
      \end{cvitems}
    }

  \cventry
    {Analyste de données}
    {Videns Analytics}
    {Montréal, Canada}
    {Avril 2020 - Juin 2021}
    {
      \begin{cvitems}
        \item{Pilotage de \textbf{missions de conseil en analytique pour l'assurance}, du nettoyage et de l'exploration des données aux modèles prédictifs et rapports automatisés pour appuyer les décisions des clients.}
        \item{Conception et animation d'un \textbf{programme de formation en Python, statistiques et apprentissage automatique}, repris par l'UQAC pour un cours universitaire. Promotion de stagiaire à analyste à temps plein.}
      \end{cvitems}
    }
\end{cventries}
